\section{Introduction}
\label{sec:intro}

Multi-site statistical analysis is increasingly constrained by data that
cannot be centralised: electronic health records behind institutional
firewalls, trials governed by separate data-use agreements, claims
databases that no single sponsor controls. Federated learning
\citep{mcmahan2017,kairouz2021} addresses this by exchanging model
parameters rather than records. A central server broadcasts a current
parameter vector, each site updates it on its own data, and the server
aggregates the returned vectors.

Let $E$ denote the number of local passes---here, full coordinate-descent
(CD) epochs---each site performs between aggregation rounds. Larger $E$
can reduce communication, but also allows local iterates to move further
toward different site-specific optima, producing client drift
\citep{li2020convergence,karimireddy2020}. For the Lasso
\citep{tibshirani1996}, both local fitting and server aggregation can affect
which coefficients are exactly zero. Consequently, objective values,
prediction error and agreement of selected variables can respond
differently to a change in $E$.

This article studies the choice of $E$ for a synchronous federated
coordinate-descent Lasso in which each site tunes its own penalty on a
local validation split and the server aggregates by sample-size weight.
The question is how to interpret changes in its estimated active set as
$E$ varies. Such changes can reflect local optimisation, averaging of
different sparse vectors, penalty selection and the numerical threshold
used to declare a coefficient active. We examine these components through
an orthogonal calculation and controlled computational comparisons.

\paragraph{Contributions.} Our contributions are as follows.

\begin{enumerate}
\item \textbf{An exact orthogonal calculation.} Under orthogonal local
designs, the CD iteration reaches
$\sum_jw_j\soft(z_j,\lambda_j)$ after one round, independently of $E$
(Proposition~\ref{prop:orth}). Except for exact cancellations, its support
is the union of the local supports and contains the pooled Lasso support.
This is a concrete calculation for the CD protocol, rather than a new
discovery that averaging may enlarge sparse supports.

\item \textbf{False positives under a fixed penalty.}
Corollary~\ref{cor:fp} gives an explicit null-selection probability under
independent Gaussian site statistics and a fixed common penalty. For
balanced partitions of a fixed sample, this probability strictly increases
with the number of sites. These qualifications are needed; the result
does not establish the same ordering for locally tuned penalties.

\item \textbf{A common objective for comparing solvers.} Sample-size
weights make the weighted local objective equal to the pooled Lasso
objective at penalty $\bar\lambda$ (Lemma~\ref{lem:pooled}). The averaged
fixed-point condition need not imply stationarity of that objective
(Proposition~\ref{prop:fixed}). Consensus ADMM \citep{boyd2011} supplies
an optimisation benchmark at the same penalty, separating objective
accuracy from support recovery.

\item \textbf{Separating tuning from aggregation.} We compare
server thresholding (\textsc{FedAvg-ST}), per-round proximal
aggregation (\textsc{FedAvg-P}), and consensus ADMM with pooled and local
baselines. Minimum-validation-error and one-standard-error variants
expose the effect of the selection rule. A communication-budget comparison
with federated dual averaging places the CD procedures alongside a
method designed for composite objectives.

\item \textbf{Communication and local work are separate costs.} We report
both transmitted bytes and local computation, including additional
federated training paths required for tuning. Increasing $E$ can save communication while
increasing total local work, so the two measures answer different
deployment questions.
\end{enumerate}

\paragraph{Related work.} The local-epoch trade-off is well studied for
smooth objectives through analyses of local SGD
\citep{stich2019,li2020convergence} and drift corrections such as FedProx
\citep{li2020fedprox} and SCAFFOLD \citep{karimireddy2020}.
Composite and proximal federated methods address nonsmooth regularisation
\citep{tran2021}. In the statistics literature, communication-efficient
distributed sparse estimation includes one-shot and surrogate-likelihood
constructions \citep{lee2017,battey2018,jordan2019} and distributed
optimisation with consensus \citep{boyd2011,duchi2012,wang2017}.

Closest to this article is the federated composite optimisation
literature. \citet{yuan2021fco} identify the ``curse of primal averaging'',
introduce FedDualAvg, and evaluate sparse recovery using precision,
recall, sparsity density and $F_1$. \citet{bao2022} develop faster
dual-averaging methods and statistical recovery guarantees, and explicitly
report support-recovery $F_1$ for sparse linear regression. Thus both the
averaging mechanism and variable selection as an evaluation target have
established precedents.

Our contribution is narrower: an explicit calculation for local cyclic
coordinate descent, a fixed-penalty null-selection formula, and a
computational assessment of local epochs under stated tuning and cost
conventions. The analysis complements established support-recovery and
stability-selection work \citep{zhaoyu2006,meinshausen2010}; it does not
provide a new general support-recovery guarantee. The simulation and data
illustration examine when the behaviour of this particular protocol
persists beyond its solvable orthogonal case.

The remainder of the article is organised as follows.
Section~\ref{sec:setup} fixes notation and states the algorithm.
Section~\ref{sec:theory} contains the analytical results.
Section~\ref{sec:method} introduces \textsc{FedAvg-ST} and the comparators.
Sections~\ref{sec:sim} and~\ref{sec:results} describe and report the Monte
Carlo study, Section~\ref{sec:budget} compares fixed training budgets, Sections~\ref{sec:realdata} and~\ref{sec:publicdata} present synthetic and public-data illustrations, and
Section~\ref{sec:discussion} discusses limitations.
